% !TeX root = ../教材题库_全部条目_转写工作稿.tex
% Source: XJ-B1-C05-EX-P189-Y0403
% 第5章 / 5.4 函数 y=A sin(ωx+φ) 的图象与性质 / 教材 p.189、190 / PDF 第196、197页
% OCR未校对：公式与文字均可能有误；未完成数学TeX转写；禁止直接发给学生
% 原题及相邻原解截图（核对用）：../题目原图/XJ-B1-C05-EX-P189-Y0403-01.webp, ../题目原图/XJ-B1-C05-EX-P189-Y0403-02.webp
% 本文件只在转写完成后作为题干使用；OCR 条目尚不是可靠题干。
\begin{bankitem}{5.4 例题 例题1}{第5章 / 5.4 函数 y=A sin(\ensuremath{\omega}x+\ensuremath{\varphi}) 的图象与性质 / 教材 p.189、190 / PDF 第196、197页}{XJ-B1-C05-EX-P189-Y0403}{OCR未校对：公式与文字均可能有误；未完成数学TeX转写；禁止直接发给学生}
\begin{ocrdraft}
[教材p.189 / PDF第196页]
\par
0二二]在同一平面直角坐标系中画出3-艹工，y=2sin产3-
\par
1
\par
Sln工
\par
在[一
\par
2
\par
丌]上的图象，观察它们之间的关系，并说出这三个函数的周期、最大值、最小值、
\par
值域之间的关系．
\par
解函数y=sin3，=2sin的周期都是2丌，在[
\par
一丌，丌]上分别求出
\par
这三个函数的图象上的五个关键点，并作出它们在一个周期内的简图，如图5．4一4．
\par
y=2s1nx
\par
2
\par
0
\par
2
\par
这三个数的奇偶
\par
性、单调区间是一样的吗？
\par
图5．4一4
\par
观察图5．4一4，可以发现：3，一2sinx的图象可以由3，=sin工的图象上每一点
\par
0，sinx）的横坐标不变、纵坐标乘2（到轴的距离放大到原来的2倍）得到，因而
\par
3，一的周期仍是2丌，最大值和最小值分别变为2和一2，值域变成[一2，2]，
\par
也就是说，“振动幅度”扩大到3，=sin的2倍．
\par
类似地，观察图5．4一4，可以发现：y=—sin的图象可以由y=sinx的图象
\par
1
\par
2
\par
，x）的横坐标不变、纵坐标乘（到，轴的距离缩短到原来的）得
\par
上每一点0
\par
到，因而的周期仍是2丌，最大值和最小值分别变为一和一
\par
，值域变为
\par
砉〕，“振
\par
1
\par
幅度”缩小为3，一的
\par
一般地，对任意A>O且，A1，函数y=Asin的图象可以由3，一sinx的图象
\par
上每一点的横坐标不变、纵坐标乘A得到．y=Asinx的周期仍是2丌，值域为[一A，
\par
[教材p.190 / PDF第197页]
\par
A)，最大值和最小值分别为，A和—A.
\par
0
\par
参数>0）对一-图象的影
\end{ocrdraft}
\end{bankitem}
